% ARCHIVO GENERADO por scripts/generar_capitulos.py. No editar a mano.
\section{Clases de la iteración 3}


\subsubsection{Dominio (CL-05)}

\diagrama{uml/CL-05_Dominio_It3.png}{CL-05}{fig:CL-05_Dominio_It3}{1.0}

\begin{itemize}
  \item \texttt{Pago} amplía sus atributos según el modelo v1.4 y agrega \texttt{valor\_descuento()} y \texttt{corregible\_por(usuario, momento)} (RN-44).
  \item \texttt{Descuento} (nueva) con \texttt{valor\_sobre(bruto)}: redondeo hacia abajo (DM-16).
  \item \texttt{CorreoConfirmacion} (nueva).
  \item \texttt{CalendarioCobro}: funciones puras del calendario de cobro (sección 5.1).
  \item \texttt{Propietario} y \texttt{Vehiculo} agregan el motivo de inactivación.
\end{itemize}

\subsubsection{Servicios (CL-06)}

\diagrama{uml/CL-06_Servicios_It3.png}{CL-06}{fig:CL-06_Servicios_It3}{1.0}

\begin{longtable}{>{\raggedright\arraybackslash}p{4.2cm}>{\raggedright\arraybackslash}p{10.2cm}}
\toprule
\textbf{Clase} & \textbf{Responsabilidad}\\
\midrule\endhead
\texttt{ServicioVigencia} (ampliado) & Solo cuenta pagos REGISTRADOS. Día de cobro (DM-12), estado de la mensualidad y meses pendientes. Las funciones \texttt{ultimo\_pago} y \texttt{esta\_vigente} conservan su nombre, así la entrada, los avisos y el mapa siguen funcionando.\\
\texttt{ServicioPagos} & \texttt{preparar} (resumen sin guardar, para M-26), \texttt{registrar} y \texttt{corregir}, cada uno en una transacción.\\
\texttt{ServicioTarifas} & Registrar tarifas y obtener la vigente en una fecha.\\
\texttt{ServicioDescuentos} & Crear, habilitar y deshabilitar descuentos.\\
\texttt{ServicioConfirmacion} & Enviar el correo simulado con el \texttt{ServicioCorreo} de \texttt{core} y guardarlo.\\
\texttt{ConsultasPagos} & Historiales, factura, por vencer e ingresos (solo lectura).\\
\texttt{ServicioPropietarios}, \texttt{ServicioVehiculos} (ampliados) & \texttt{cambiar\_estado} recibe el motivo; \texttt{reactivar\_por\_pago} solo reactiva si el motivo es FALTA\_PAGO.\\
\bottomrule
\end{longtable}



